\documentclass[11pt]{ctexart}
\usepackage[a4paper,top=2.2cm,bottom=2.2cm,left=2.2cm,right=2.2cm]{geometry}
\usepackage{booktabs}
\usepackage{longtable}
\usepackage{array}
\usepackage{enumitem}
\usepackage{xcolor}
\usepackage{amssymb}
\usepackage[hidelinks]{hyperref}
\setlist{nosep, leftmargin=1.6em}
\emergencystretch=1.5em
\setlength{\tabcolsep}{4pt}
\renewcommand{\arraystretch}{0.92}

\hypersetup{
  pdftitle={SmartFarmAgent 智慧农业智能体系统——项目方案},
  pdfsubject={项目方案（导师提交版，整体重写 v2.0）},
  pdfauthor={}
}

\begin{document}

\begin{center}
  {\LARGE\bfseries SmartFarmAgent 智慧农业智能体系统}\\[4pt]
  {\large 项目方案}
\end{center}

\vspace{0.5em}

\section{项目背景与痛点}
\begin{itemize}
  \item 行业痛点：传统种养管理依赖人工巡检与经验调控，存在病害发现滞后、水肥浪费、环境调控不及时等问题；企业级智慧农业平台功能完整但重软件集成、成本高，市面常见产品以阈值规则控制为主，大模型应用停留在“知识问答”层，均未把智能决策闭环到设备执行。
  \item 竞赛方向：农业农村部 2026 年智慧农业创新大赛六大赛道（智能农机作业控制、无人机巡田、激光除草机器人、设施小番茄采摘机器人、家禽巡检机器人、鱼群智能检测装备）集中在“智能体 / IoT + 机器人”方向，强调感知、决策、执行一体化落地。
  \item 技术窗口：农业开源大模型与 RAG 框架成熟可用（司农、RAGFlow、Dify 等）；ESP32 单芯片即可完成数据采集、拍照与控制，Node.js + Vue 平台生态成熟，硬件与软件成本都低，适合以低成本原型参赛。
\end{itemize}

\section{与 2026 智慧农业创新大赛的结合点}

本项目以低成本硬件原型实现“感知—决策—执行”一体化闭环，与大赛各赛道对设备级智能化落地的要求契合。结合点如下：

\begin{longtable}{p{3.6cm}p{5.0cm}p{6.6cm}}
  \toprule
  大赛赛道 & 本项目对应能力 & 结合说明 \\
  \midrule
  智能农机作业控制 & ESP32 设备控制链路（继电器/泵/灯/喷药） & 设备级“指令—执行—回传”闭环可迁移到农机作业控制单元 \\
  无人机巡田解决方案 & 视觉识别 + 智能体决策链路 & 巡田影像的“识别—分析—处方”流程与本项目视觉+决策链路同构 \\
  设施采摘、巡检机器人 & ESP32-CAM 视觉 + 执行器联动 & 视觉识别触发设备动作的架构可复用到采摘、巡检场景 \\
  家禽巡检、鱼群检测 & 感知上云 + 异常告警 + 自动处置 & 传感采集与云端判别的模式一致，更换传感器与模型即可适配 \\
  \bottomrule
\end{longtable}

综合来看，本项目与「设施小番茄采摘机器人」赛道契合度最高：两者同属设施农业场景，视觉识别、智能体决策与执行器联动的链路一致，本项目的 ESP32-CAM 视觉与控制链可直接复用，只需把浇水、喷药等执行器替换为采摘装置；与「智能农机作业控制」赛道则在设备级控制闭环上同构。其余赛道需要更换感知与执行模块，但“感知—决策—执行”的通用底座可以复用。

\textbf{定位：}以 ESP32 感知执行端 + Node.js/Vue 平台 + 云端大模型智能体，做一个面向智慧农业竞赛与科研的低成本“闭环能力原型”，聚焦大模型决策到设备动作的完整链路，可针对具体赛道快速适配。

\section{系统总体架构}
\begin{longtable}{p{1.2cm}p{2.2cm}p{8.8cm}p{3.0cm}}
  \toprule
  层级 & 名称 & 组成与职责 & 关键技术 \\
  \midrule
  L6 & 展示层 & Vue 3 Web 大屏 / H5：实时曲线、设备控制、农技问答、视觉告警 & Vue、ECharts、WebSocket \\
  L5 & 智能决策层 & 大模型智能体（LLM+RAG）→ JSON 设备指令 + 农技问答（核心创新）；视觉服务（YOLO 病害/生长分析） & LLM、RAG、视觉 \\
  L4 & 云端平台层 & Node.js 后端：MQTT 接入（mqtt.js）→ 时序存储 → REST/WebSocket API；智能体编排 & Node.js、MQTT、时序存储 \\
  L3 & 边缘层 & ESP32-CAM 定时拍照上传；断网规则兜底（本地优先） & 边缘计算、定时任务 \\
  L2 & 通信层 & ESP32 内置 Wi-Fi：MQTT（JSON 上行/下行）+ HTTP 图片上传 & Wi-Fi、MQTT \\
  L1 & 感知执行层 & ESP32 单芯片：传感器采集 + 继电器/泵/灯/喷药 + OLED + microSD 日志 & GPIO、ADC、I2C、PWM \\
  \bottomrule
\end{longtable}
安全原则：本地阈值规则永远优先于云端指令（断网自持、误操作兜底），云端智能体指令走白名单。

\section{功能模块}
\begin{small}
\begin{longtable}{p{4.4cm}p{5.4cm}p{1.0cm}p{4.4cm}}
  \toprule
  功能 & 实现方案 & 难度 & 可行性依据 \\
  \midrule
  自动浇水/施肥 & 土壤湿度+光照综合判断；继电器驱动水泵/蠕动泵 & 中 & 已有土壤传感器+继电器；ESP32 整机案例已验证 \\
  环境自动调控 & DHT11/BH1750 → 风扇/加热/补光灯；PWM 调光（PID 可选） & 中 & 传感器已有；PWM 知识可迁移 \\
  病虫害检测+给药 & ESP32-CAM 定时拍照上传 → YOLOv8（先云后边）→ 置信度达标 → 喷药泵 & 高 & PlantDoc/PlantVillage 公开数据集 + ESP32-CAM/YOLO 全链路先例 \\
  生长状况分析 & 定距定时拍照 → 叶片数/株高/颜色统计 → 大模型生成周报 & 高 & 视觉测量工具包开源可用 \\
  农技知识问答 & 种植/病虫害/用药知识库 → RAG 检索增强 → 结构化方案+来源标注 & 中 & 司农/DeepSeek + RAGFlow 已核验，Node.js 可调其 HTTP API \\
  智能体决策 & Node.js 编排：传感器+视觉+历史 → 大模型（司农/DeepSeek）→ 指令+解释 & 高 & LLM 兼容 API + MQTT 下行链路已验证 \\
  数据中台/展示 & MQTT 上行 → Node.js 时序存储 → Vue 大屏；microSD 断网日志 & 中 & Express/mqtt.js/ECharts 成熟；EMQX 已核验 \\
  \bottomrule
\end{longtable}
\end{small}

\section{数据链路}
\begin{longtable}{p{1.5cm}p{11.0cm}p{2.6cm}}
  \toprule
  环节 & 内容 & 输出 \\
  \midrule
  感知 & BH1750/DHT11/土壤 → ESP32 采集打包（JSON）；ESP32-CAM 定时拍照 & 数据 / 图片 \\
  通信 & ESP32 Wi-Fi：MQTT 上行情报 + HTTP 图片上传；断网时 microSD 本地日志 & 上行数据 / 图片 \\
  云端 & Node.js 后端（mqtt.js）订阅 → 时序存储 → REST/WebSocket API & 时序数据 \\
  AI & 视觉服务（YOLO 病害检测/生长分析）→ 结果入事件队列；智能体（LLM+RAG）综合传感器+视觉+历史，输出设备指令与农技方案 & 决策 JSON / 问答方案 \\
  执行 & MQTT 下行指令 → ESP32 解析（白名单）→ 继电器/泵/灯/喷药 & 设备动作 \\
  闭环 & 执行状态回传 → 数据入库 → Vue 大屏展示 + 智能体下次决策依据 & 反馈数据 \\
  \bottomrule
\end{longtable}

\section{技术可行性}
\begin{enumerate}
  \item 硬件与嵌入式：ESP32 生态成熟（arduino-esp32 官方核心，17k+ 星，持续更新）；ESP32-CAM 单板集成摄像头、Wi-Fi 与 microSD，外接传感器/继电器案例丰富；现有 I2C、SPI、串口、PWM 知识可直接迁移。
  \item 视觉检测：PlantDoc（427 星）/ PlantVillage（5 万+ 图、38 类）公开数据集可用，YOLOv8 框架成熟（Ultralytics 60k 星），已有温室病害检测与 ESP32-CAM 全链路检测先例。
  \item 智能体链路：司农（8B/32B）与 AgriAgent（125 星）均已开源；RAGFlow/Dify 提供 HTTP API，Node.js 后端可直接调用，智能体到设备路径可行。
  \item 平台：Express（69k 星）、MQTT.js（9k 星）、Vue 3（54k 星）、ECharts（67k 星）均已核验；EMQX（16.5k 星）与 OneNET 支持 ESP32 直连，免费可用。
  \item 已知约束：ESP32-CAM 引脚少、供电敏感，需稳压供电与 I/O 规划；初期视觉推理走云端；司农 32B 本地部署显存要求高，初期用 8B 或 DeepSeek API。
\end{enumerate}

\section{里程碑计划}
\begin{longtable}{p{2.2cm}p{1.7cm}p{6.5cm}p{4.0cm}}
  \toprule
  阶段 & 时间 & 内容 & 验收点 \\
  \midrule
  M0 开题 & 第 1—2 周 & 链路定稿（ESP32 + Node.js + Vue）、软件环境搭建 & 开题通过 \\
  M1 感知-控制闭环 & 第 3—6 周 & ESP32 传感器采集→OLED→继电器→MQTT（本地 EMQX）→Node.js→Vue 控制页 & 局域网内数据正确、远程手动可控 \\
  M2 上云+拍照+日志 & 第 7—10 周 & ESP32-CAM 定时拍照上传；MQTT 上云；Node.js 时序存储+Vue 大屏；microSD 断网补传 & 云平台实时曲线、图片可查、断网日志完整可补传 \\
  M3 视觉检测 & 第 11—18 周 & 公开集+自采数据训练 YOLO（先 5 类常见病）→Node.js 调视觉服务→触发喷药 & 检测准确率 ≥85\%，喷药动作联动 \\
  M4 智能体决策 & 第 19—23 周 & Node.js 编排智能体+RAG 知识库（设备决策+农技问答）→ 指令闭环 → Vue 对话 & 自然语言指令自动执行，农技问答可溯源，误操作有兜底 \\
  M5 作品化 & 第 24 周 & 外壳/文档/演示视频/说明书/答辩材料 & 达到参赛交付标准 \\
  \bottomrule
\end{longtable}

\section{创新点与预期成果}
\begin{enumerate}
  \item 智能体闭环决策：大模型直接输出设备动作指令，不只停留在问答层。
  \item 边缘-云协同：断网时本地规则独立运行，联网后云端智能体接管；本地规则永远优先，误操作有兜底。
  \item 低成本单芯片可复制：ESP32 单芯片完成采集、拍照、控制与联网，链路最短、硬件总成本低，整套闭环可复现。
  \item 全周期数字档案：从种子到收获持续记录传感器与视觉数据，形成可追溯的生长档案。
\end{enumerate}

\section{风险与应对}
\begin{longtable}{p{5.2cm}p{10.0cm}}
  \toprule
  风险 & 应对 \\
  \midrule
  视觉识别数据不足/效果差 & 公开数据集预训练 + 自采微调；范围先缩到 3—5 类常见病害 \\
  ESP32-CAM 算力不足 & 初期视觉推理放云端，边缘只拍照上传；后期再试轻量模型 \\
  ESP32-CAM 引脚少/供电敏感 & 稳压供电；继电器组用 I/O 扩展（如 PCF8574/74HC595）或精简执行器数量 \\
  智能体误操作 & 本地规则优先、指令白名单、关键动作人工确认模式 \\
  Wi-Fi 不稳定 & 本地阈值模式 + microSD 日志补传 \\
  时间不足 & 裁剪顺序：合并施肥/给药、生长分析降级为拍照+周报、Vue 先做单页大屏 \\
  \bottomrule
\end{longtable}

\section{主要开源参考}
\begin{itemize}
  \item \url{https://github.com/langgenius/dify}（智能体/工作流平台，151k 星）
  \item \url{https://github.com/infiniflow/ragflow}（RAG 引擎，87k 星）
  \item \url{https://github.com/ultralytics/ultralytics}（YOLO 框架，60k 星）
  \item \url{https://github.com/pratikkayal/PlantDoc-Dataset}（病害数据集，427 星）
  \item PlantVillage（\url{https://www.kaggle.com/datasets/abdallahalidev/plantvillage-dataset}，5 万+ 图、38 类）
  \item 司农（南京农业大学农业大语言模型，8B/32B）
  \item \url{https://github.com/zhiweihu1103/AgriAgent}（农业多模态大模型，125 星）
  \item \url{https://github.com/expressjs/express}（69k 星）/ \url{https://github.com/mqttjs/MQTT.js}（9k 星）（Node.js 后端）
  \item \url{https://github.com/vuejs/core}（54k 星）/ \url{https://github.com/apache/echarts}（67k 星）（前端展示）
  \item \url{https://github.com/emqx/emqx}（16.5k 星）（MQTT Broker）
\end{itemize}

\end{document}